\begin{lemma}
	\label{lem:cofactor-same}
	If \( \bA = (a_{ij}) \) is a square matrix of \( n\times n \) dimension, s.t.
	\begin{itemize}
		\item all the row sums are zero.
		\item all the column sums are zero.
	\end{itemize}
	then:
	\begin{enumerate}
		\item for \( 1\leq r \leq n \), have:
		      \[
			      A_{r,1} = A_{r,2} = \cdots = A_{r,n}
		      \]
		\item for \( 1\leq c\leq r \)
		      \[
			      A_{1,c} = A_{2,c} = \cdots = A_{n,c}
		      \]
	\end{enumerate}
	Namely all its cofactors are the same.
\end{lemma}
\begin{proof}
	Fix \( r \), let \( \bA' \) be the matrix that has \( a'_{r,1} = a_{r,1} + 1 \), which bump the row sum of row \( r \) to be \( 1 \) instead of \( 0 \), and the other entries remains the same as \( \bA \). Choose a \( k \in \br{1,n} \). Let \( \bA'' \) be the matrix with the \( k \)-th column equal to the sum of all columns of \( \bA' \), and the other columns the same as \( \bA' \). As the column sum will be \( 0 \) except for row \( r \), the $k$-column of \( \bA'' \) will admit \( 1 \) at \( (r,k) \) and \( 0 \) at the rest of the entries. We have following observation:
	\begin{itemize}
		\item \( \det \bA' = \det\bA'' \). Since \( \bA'' \) is obtained by doing elementary operation from \( \bA' \).
		\item \( \det \bA'' = A''_{r,k} \) by the fact that the \( k \)-th column only admits \( 1 \) at \( r \)-th row and \( 0 \) elsewhere.
		\item \( \det \bA + A_{r,1} = \det\bA' \) by definition of determinant, and notice that \( \det\bA = 0 \implies A_{r,1} = A''_{r,k}\), and by construction of \( \bA'' \) see that \( A''_{r,k} = A_{r,k} \).
	\end{itemize}
	becase \( \bA' \) is independent of \( k \), \( A_{r,k} \) is independent of \( k \).
\end{proof}

\begin{definition}[\textbf{\( d \)-Regular Graph}]
	A graph is \underline{\( d \)-regular} if all vertices have degree to be \( d \).
\end{definition}

\begin{proposition}
	If \( G \) is \( d \)-regular on \( n \) vertices, the \( \LL(G) = d \cdot \II_{n} - \bA(G) \). In particular, if \( G \) has eigenvalues to be \( \lambda_1, \ldots, \lambda_n \) with \( \lambda_n = d \), then \( \LL(G) \) has eigenvalues \( d-\lambda_1, \ldots, d-\lambda_{n-1}, d-\lambda_n = 0 \). Moreover:
	\[
		\kappa(G) = \frac{1}{n} \prod_{i=1}^{n-1} (d-\lambda_i)
	\]
\end{proposition}

\begin{proof}
	Since \( \LL(G) \) has zero row and column sums, cofactors of \( \LL(G) \) are equal by \textbf{Lemma} \ref{lem:cofactor-same}. Compute the characteristic polynomial of \( \LL(G) \):
	\[
		\begin{aligned}
			\det(\LL(G) - t\cdot \II_n) & = \det(\LL(G)) - t(L_1 + \cdots + L_{n,n}) + \text{higher degree terms}\quad (\text{results in linear algebra}) \\
			                            & = \det(\LL(G)) - nt L_{1,1} + \text{higher degree terms}                                                        \\
			                            & = -ntL_{1,1} + \text{higher degree terms}
		\end{aligned}
	\]
	On the other hand,
	\[
		\det(\LL(G) - t\II_n) = \prod_{i=1}^n (\underbrace{\lambda'_i}_{\text{e.v. of } \LL(G)} - t) = -t\underbrace{(d-\lambda_1) \cdots (d-\lambda_{n-1})}_{\text{eigenvalues of } \LL(G) = d \cdot \II_n - \bA(G)} + \text{higher degree terms}
	\]
	thus
	\[
		\kappa(G) = L_{i,i} = \frac{1}{n} \prod_{i=1}^{n-1} (d-\lambda_i)
	\]
\end{proof}

\begin{corollary}
	Let \( G \) be a connected loopless graph with \( p \) vertices. Suppose that the eigenvalues of \( \LL(G) \) are given by \( \lambda_1, \ldots, \lambda_{n-1}, \lambda_n = 0 \), then:
	\[
		\kappa(G) = \frac{1}{n}\prod_{i=1}^{n-1}\lambda_i
	\]
\end{corollary}

\begin{eg}
	\leavevmode
	\begin{enumerate}
		\item Consider \( C_d \) be the cube graph with \( d \) dimensions, total vertices to be \( 2^d \), have:
		      \begin{itemize}
			      \item \( \bA(C_d) \) has \( d-2i \) as eigenvalues with multiplicity to be \( \binom{d}{i} \) for \( i \in \br{1,d} \).
			      \item \( \LL(C_d) \) has eigenvalues \( 2i \) with multiplicity being \( \binom{d}{i} \) for \( i\in \br{1,d} \).
			      \item Thus:
			            \[
				            \kappa(C_d)  =\frac{1}{2^d} \prod_{i=1}^d (2i)^{\binom{d}{i}}
			            \]
		      \end{itemize}
		\item Consider \( K_p \) be the complete graph with \( p \) vertices.
		      \begin{itemize}
			      \item \( \bA(K_p) \) has eigenvalues being \( p-1, -1, \ldots, -1 \).
			      \item \( \LL(K_p) \) has eigenavalues being \( 0, p, \ldots, p \).
			      \item Thus:
			            \[
				            \kappa(K_p) = \frac{1}{p} \cdot (p^{p-1}) = p^{p-2}
			            \]
		      \end{itemize}
	\end{enumerate}
\end{eg}

\section{Eulerian Tours}

\begin{definition}[\textbf{Binary de Bruijn}]
	Given a positive integer \( k \), a \underline{binary de Bruijn sequence of degree \( k \)} is a sequence of \( 0,1 \) that is periodic of period \( 2^k \) and that has the list of all contiguous subsequence to be exactly all binary strings of length \( k \).
\end{definition}

Namely we scan through the string, and any kinds of sequence consisting that alphabet with length \( k \) should only appear once.

\begin{eg}
	\leavevmode
	\begin{itemize}
		\item When \( k=2 \), the following forms a binary de Bruijn sequence of degree \( 2 \).
		      \[
			      1,1,0,0,1
		      \]
		\item When \( k=3 \), the following forms a binary de Bruijn sequence of degree \( 3 \).
		      \[
			      0,0,0,1,1,1,0,1,0,0,0
		      \]
		      which attains \( 000 \), \( 001 \), \( 011 \), \( 111 \), \( 101 \), \( 010 \), \( 100 \).
	\end{itemize}
\end{eg}

\begin{definition}[\textbf{Directed graph}]
	A \underline{directed graph or a digraph} \( D \) is \( (V,E, \vp) \) where \( \vp: E \to V \times V \) which maps to ordered pair of vertices.
\end{definition}

\begin{definition}[\textbf{Eulerian Tour}]
	An \underline{Eulerian tour} in \( G \) is a closed walk that visit all vertices and very edge exactly \( 1 \) times. Note that vertices can be visited many times.
\end{definition}

